\documentclass[12pt,a4paper]{article}

% -------------------------------------------------
% Sprache / Kodierung
% -------------------------------------------------
\usepackage[T1]{fontenc}
\usepackage[utf8]{inputenc}
\usepackage[ngerman]{babel}
\usepackage{lmodern}
\usepackage{microtype}

% -------------------------------------------------
% Mathematik
% -------------------------------------------------
\usepackage{amsmath,amssymb,amsthm,mathtools}

% -------------------------------------------------
% Layout / Grafik / Farben
% -------------------------------------------------
\usepackage[a4paper,margin=2.3cm]{geometry}
\usepackage{booktabs}
\usepackage{xcolor}
\usepackage[hidelinks,unicode]{hyperref}
\pdfstringdefDisableCommands{%
  \def\ü{ue}\def\ö{oe}\def\ä{ae}\def\Ü{Ue}\def\Ö{Oe}\def\Ä{Ae}\def\ss{ss}%
}
\usepackage[most]{tcolorbox}
\usepackage{enumitem}
\usepackage{array}

% -------------------------------------------------
% Farben (konsistent mit Vorgängerdokumenten)
% -------------------------------------------------
\definecolor{lückenrot}{RGB}{180,40,40}
\definecolor{bewiesengrün}{RGB}{30,100,50}
\definecolor{warnorange}{RGB}{200,100,10}
\definecolor{hintergrundblau}{RGB}{235,242,250}
\definecolor{hintergrundgelb}{RGB}{255,252,225}
\definecolor{adischviolett}{RGB}{90,20,140}
\definecolor{fehlerrot}{RGB}{200,0,0}
\definecolor{konjviolett}{RGB}{80,0,130}
\definecolor{e8grau}{RGB}{55,55,90}
\definecolor{korrekturorange}{RGB}{200,100,10}
\definecolor{bernoulliblau}{RGB}{0,70,130}
\definecolor{quaterniongrün}{RGB}{0,100,60}

% -------------------------------------------------
% tcolorbox-Stile
% -------------------------------------------------
\tcbset{
  lücke/.style={colback=red!8, colframe=lückenrot, fonttitle=\bfseries,
    title={Offene Lücke}},
  ergebnis/.style={colback=green!7, colframe=bewiesengrün,
    fonttitle=\bfseries, title={Bewiesenes Ergebnis}},
  warnung/.style={colback=orange!10, colframe=warnorange,
    fonttitle=\bfseries, title={Achtung / Heuristik}},
  adisch/.style={colback=violet!6, colframe=adischviolett,
    fonttitle=\bfseries, title={2-adische Perspektive}},
  kernidee/.style={colback=hintergrundblau, colframe=blue!60!black,
    fonttitle=\bfseries, title={Kernidee}},
  konj/.style={colback=violet!5, colframe=konjviolett, fonttitle=\bfseries,
    title={Konjektur (unbewiesen)}},
  lte/.style={colback=yellow!8, colframe=e8grau, fonttitle=\bfseries,
    title={Modellrechnung}},
  hauptsatz/.style={colback=green!10, colframe=bewiesengrün!80!black,
    fonttitle=\bfseries\large, title={Stärkster beweisbarer Satz}},
  korrektur/.style={colback=yellow!5, colframe=orange!70!black,
    fonttitle=\bfseries},
  bernoulli/.style={colback=blue!5, colframe=bernoulliblau,
    fonttitle=\bfseries},
  quaternion/.style={colback=green!5, colframe=quaterniongrün,
    fonttitle=\bfseries, title={Quaternionen-Perspektive}},
}

% -------------------------------------------------
% Theorem-Umgebungen
% -------------------------------------------------
\newtheorem{definition}{Definition}[section]
\newtheorem{proposition}[definition]{Proposition}
\newtheorem{theorem}[definition]{Theorem}
\newtheorem{lemma}[definition]{Lemma}
\newtheorem{korollar}[definition]{Korollar}
\newtheorem{bemerkung}[definition]{Bemerkung}
\newtheorem{hypothese}[definition]{Hypothese}
\newtheorem{satz}[definition]{Satz}

\newenvironment{beweis}{\begin{proof}[Beweis]}{\end{proof}}

% -------------------------------------------------
% Makros
% -------------------------------------------------
\newcommand{\vii}{\nu_2}
\newcommand{\U}{\mathcal{U}}
\newcommand{\N}{\mathbb{N}}
\newcommand{\Z}{\mathbb{Z}}
\newcommand{\R}{\mathbb{R}}
\newcommand{\C}{\mathbb{C}}
\newcommand{\E}{\mathbb{E}}
\newcommand{\Q}{\mathbb{Q}}
\newcommand{\HH}{\mathbb{H}}
\newcommand{\Ztwo}{\mathbb{Z}_2}
\newcommand{\abs}[1]{\lvert#1\rvert}
\newcommand{\atwo}[1]{\lvert#1\rvert_2}
\newcommand{\logt}{\log_2}
\newcommand{\In}{\mathcal{I}}
\newcommand{\Denom}[1]{D(#1)}
\newcommand{\Nq}{N}          % Quaternionennorm
\newcommand{\Re}{\operatorname{Re}}

% -------------------------------------------------
% Dokument
% -------------------------------------------------
\title{%
  Vier-Teilnormen-Faktorisierung und Collatz-Dynamik:\\[0.3em]
  Asymmetrische Quaternionentransformation unter $3n+1$\\[0.5em]
  \large Lagrange-Vierquadratesatz, Hurwitz-Quaternionen und\\
  optimale Quaternionenwahl als Minimierungsproblem\\[0.3em]
  \large Ergänzung zur Collatz-Synthese vom 16.\ Juni 2026%
}
\author{Thomas Hoffbauer}
\date{16.\ Juni 2026}

\begin{document}
\maketitle

\begin{abstract}
Wir untersuchen einen neuen Zugang zur Collatz-Dynamik über die
\emph{Vier-Quadrate-Zerlegung} im Sinne des Lagrange-Vierquadratensatzes:
Jede natürliche Zahl $n$ besitzt eine Darstellung $n = a^2+b^2+c^2+d^2$
als Norm eines Quaternions $q = a+bi+cj+dk \in \HH$.
Der Collatz-Schritt $n \mapsto 3n+1$ wirkt auf der Quaternionen-Ebene
als \emph{asymmetrische Transformation}: Die Abbildung $q \mapsto 3q+1$
transformiert die vier Teilnormen $(N_1,N_2,N_3,N_4) = (a^2,b^2,c^2,d^2)$
mit unterschiedlichen Korrekturtermen:
\[
  \Nq(3q+1) \;=\; 9n + 6\,\Re(q) + 1 \;=\; 9n + 6a + 1.
\]
Nur die erste Teilnorm $N_1 = a^2$ erhält den Korrekturfaktor $6a+1$,
während $N_2, N_3, N_4$ uniform um Faktor~$9$ wachsen. Diese \emph{ausgezeichnete
Richtung} (die Realteil-Achse) trägt den gesamten Collatz-Fehlerterm.

Das führt zu einem natürlichen Minimierungsproblem: Wähle die Zerlegung so,
dass $a$ minimal (negativ) ist. Für Quadratzahlen $n = m^2$ mit $a = -m$
ergibt sich die bemerkenswerte Formel $\Nq(3q+1) = (3m-1)^2$ --- eine
Quadratzahl. Ferner untersuchen wir Korrelationen zwischen der EABC-Klassen-Struktur
und den Paritätsmustern der Vier-Quadrate-Koeffizienten sowie einen Ausblick
auf die D${}_4$-Theta-Funktion und Bernoulli-Normschalen $\In(n)$.
\end{abstract}

\tableofcontents

\bigskip
\begin{tcolorbox}[warnung, title={Einordnung}]
Die Collatz-Vermutung ist unbewiesen. Dieses Dokument präsentiert einen
\textbf{explorativen algebraischen Zugang}, der auf der Quaternionen-Darstellung
von $n$ beruht. Die Vier-Teilnormen-Zerlegung liefert interessante Strukturen,
aber noch keinen vollständigen Konvergenzbeweis.
\end{tcolorbox}

% ====================================================
\section{Vier-Quadrate-Zerlegung und Quaternionendarstellung}
\label{sec:zerlegung}
% ====================================================

\subsection{Der Satz von Lagrange}

\begin{theorem}[Lagrange, 1770]
\label{thm:lagrange}
Jede positive ganze Zahl $n \in \N$ lässt sich als Summe von vier ganzzahligen
Quadraten schreiben:
\[
  n \;=\; a^2 + b^2 + c^2 + d^2, \quad a,b,c,d \in \Z.
\]
\end{theorem}

Die Anzahl der Darstellungen ist durch den Jacobi-Vierquadratensatz gegeben:
\[
  r_4(n) \;=\; 8 \sum_{\substack{d \mid n \\ 4 \nmid d}} d.
\]
Insbesondere gilt $r_4(n) \neq 0$ für alle $n \geq 1$, und für Primzahlen $p$
ist $r_4(p) = 8(p+1) > 0$.

\subsection{Lipschitz- und Hurwitz-Quaternionen}

\begin{definition}[Lipschitz-Quaternionen]
Die \emph{Lipschitz-Quaternionen} sind $\mathcal{L} := \Z + \Z i + \Z j + \Z k \subset \HH$,
d.\,h.\ Quaternionen mit ganzzahligen Koeffizienten.
Für $q = a + bi + cj + dk \in \mathcal{L}$ ist die \emph{Norm}
\[
  N(q) \;:=\; q\bar{q} \;=\; a^2 + b^2 + c^2 + d^2 \;\in\; \N_0.
\]
Nach Lagrange gilt: Für jedes $n \in \N$ existiert $q \in \mathcal{L}$ mit $N(q) = n$.
\end{definition}

\begin{definition}[Hurwitz-Quaternionen]
Die \emph{Hurwitz-Quaternionen} sind
\[
  \mathcal{H} \;:=\; \Bigl\{ a + bi + cj + dk \;\Big|\;
    a,b,c,d \in \Z \;\text{ oder alle } \in \Z + \tfrac{1}{2} \Bigr\}.
\]
$\mathcal{H}$ ist ein maximaler Ordnungsring in $\HH$ und bildet das D${}_4$-Gitter.
Jede ungerade Primzahl $p$ ist Norm eines Hurwitz-Quaternions: $p = N(q_p)$ für
ein geeignetes $q_p \in \mathcal{H}$.
\end{definition}

\subsection{Die vier Teilnormen}

\begin{definition}[Vier Teilnormen]
\label{def:teilnormen}
Sei $q = a + bi + cj + dk \in \mathcal{L}$ mit $N(q) = n$. Die \emph{vier Teilnormen}
(oder Quadrantnormen) sind:
\[
  N_1 \;:=\; a^2, \quad
  N_2 \;:=\; b^2, \quad
  N_3 \;:=\; c^2, \quad
  N_4 \;:=\; d^2,
\]
mit der Zerlegungseigenschaft $N(q) = N_1 + N_2 + N_3 + N_4 = n$.
\end{definition}

\begin{tcolorbox}[kernidee]
Die Vier-Teilnormen-Faktorisierung bricht die Norm $n$ in vier unabhängige Quadratbeiträge
auf. Die Frage ist: Wie transformiert der Collatz-Schritt $n \mapsto 3n+1$ diese vier
Beiträge --- und gibt es eine ausgezeichnete Zerlegung, die die Collatz-Dynamik
besonders transparent macht?
\end{tcolorbox}

% ====================================================
\section{Collatz-Schritt als Quaternionentransformation}
\label{sec:transformation}
% ====================================================

\subsection{Hauptformel}

\begin{proposition}[Asymmetrische Quaternionentransformation]
\label{prop:hauptformel}
Sei $q = a + bi + cj + dk \in \mathcal{L}$ mit $N(q) = n$. Dann gilt:
\[
  \boxed{N(3q+1) \;=\; 9\,N(q) + 6\,\Re(q) + 1 \;=\; 9n + 6a + 1.}
\]
\end{proposition}

\begin{beweis}
Sei $q = a + bi + cj + dk$. Dann ist $3q + 1 = (3a+1) + 3bi + 3cj + 3dk$.
Die Quaternionen-Norm ist die Summe der Quadrate der Koeffizienten:
\begin{align*}
  N(3q+1) &= (3a+1)^2 + (3b)^2 + (3c)^2 + (3d)^2 \\
           &= 9a^2 + 6a + 1 + 9b^2 + 9c^2 + 9d^2 \\
           &= 9(a^2+b^2+c^2+d^2) + 6a + 1 \\
           &= 9\,N(q) + 6\,\Re(q) + 1. \qedhere
\end{align*}
\end{beweis}

\subsection{Transformation der vier Teilnormen}

Die Proposition~\ref{prop:hauptformel} lässt sich auf Teilnormen-Ebene präzisieren.
Unter $q \mapsto 3q+1$ transformieren sich die Teilnormen wie folgt:
\begin{align}
  N_1(3q+1) &= (3a+1)^2 = 9a^2 + 6a + 1 = 9N_1 + 6\sqrt{N_1} + 1, \label{eq:N1}\\
  N_2(3q+1) &= (3b)^2   = 9b^2            = 9N_2, \label{eq:N2}\\
  N_3(3q+1) &= (3c)^2   = 9c^2            = 9N_3, \label{eq:N3}\\
  N_4(3q+1) &= (3d)^2   = 9d^2            = 9N_4. \label{eq:N4}
\end{align}

\begin{tcolorbox}[quaternion]
\textbf{Kernbeobachtung: Asymmetrie der Transformation.}

Drei Teilnormen $N_2, N_3, N_4$ wachsen uniform um \emph{Faktor~$9$}.
Nur $N_1 = a^2$ (die Realteil-Teilnorm) erhält den \emph{Korrekturfaktor} $6a+1$.
Die Abbildung $q \mapsto 3q+1$ ist daher \textbf{nicht isotrop}: Sie zerstört
die Symmetrie zwischen den vier Teilnormen und definiert eine ausgezeichnete
Richtung --- die Realteil-Richtung $\Re(q)$.
\end{tcolorbox}

\begin{bemerkung}[Geometrische Interpretation]
In $\R^4 \cong \HH$ wirkt $q \mapsto 3q$ als isotrope Skalierung (alle Koordinaten
mal~$3$). Der Versatz $q \mapsto q + \frac{1}{3}$ verschiebt ausschließlich in
Realteil-Richtung. Die Collatz-Abbildung $q \mapsto 3q + 1$ ist daher eine
\emph{Kombination aus isotroper Streckung und nicht-isotroper Translation},
wobei die Translation $+1$ die Realteil-Richtung auszeichnet.
\end{bemerkung}

\subsection{Wann ist \texorpdfstring{$a$}{a} wichtig?}

Der Korrekturfaktor $6a + 1$ bestimmt den Unterschied zwischen dem Collatz-Bild
$9n+6a+1$ und dem isotropen Bild $9n+1$. Für $a > 0$ expandiert die Transformation
stärker als isotrop, für $a < 0$ schwächer.

\begin{itemize}[noitemsep]
  \item $a > 0$: $N(3q+1) = 9n + 6a + 1 > 9n + 1$ --- \emph{Überkompensation}.
  \item $a = 0$: $N(3q+1) = 9n + 1$ --- \emph{isotropes Bild}.
  \item $a < 0$: $N(3q+1) = 9n + 6a + 1 < 9n + 1$ --- \emph{Unterkompensation} (Kontraktion relativ zur isotropen Referenz).
\end{itemize}

Entscheidend für die Collatz-Dynamik ist nicht $N(3q+1)$, sondern
$N(3q+1)/2^{\nu_2(3n+1)}$ nach Division durch die maximale 2er-Potenz.
Die Frage, ob $a$ die 2-adische Valuation $\vii(3n+1)$ beeinflusst, führt
zur Verbindung mit den EABC-Klassen (Abschnitt~\ref{sec:eabc}).

% ====================================================
\section{EABC-Klassen und Vier-Quadrate-Struktur}
\label{sec:eabc}
% ====================================================

\subsection{EABC-Klassifikation}

Das EABC-Modell klassifiziert ungerade Zahlen $n$ nach $\vii(3n+1)$:
\begin{itemize}[noitemsep]
  \item \textbf{E}: $n \equiv 1 \pmod{4}$, also $\vii(3n+1) \geq 2$.
  \item \textbf{A}: $n \equiv 3 \pmod{8}$, also $\vii(3n+1) = 1$.
  \item \textbf{B}: $n \equiv 7 \pmod{16}$, also $\vii(3n+1) = 1$ (nach anderem Muster).
  \item \textbf{C}: $n \equiv 5 \pmod{8}$, also $\vii(3n+1) \geq 3$.
\end{itemize}
(Genaue Definition und Beweis in den Vorgängerdokumenten der EABC-Serie.)

\subsection{Paritätsmuster in der Vier-Quadrate-Zerlegung}

Der Jacobi-Vierquadratensatz gibt Aufschluss über die \emph{minimalen} Zerlegungen.
Ein zentrales Resultat über die Paritätsstruktur:

\begin{proposition}[Paritätsmuster nach Residuenklasse]
\label{prop:paritaet}
Sei $n$ ungerade. In jeder Vier-Quadrate-Zerlegung $n = a^2+b^2+c^2+d^2$ gilt:
\begin{enumerate}[label=(\roman*), noitemsep]
  \item $n \equiv 1 \pmod{4}$: Die Anzahl ungerader Koeffizienten unter $a,b,c,d$ ist
        gerade (0, 2 oder 4 ungerade Koeffizienten).
  \item $n \equiv 3 \pmod{4}$: Die Anzahl ungerader Koeffizienten ist ungerade (1 oder 3).
\end{enumerate}
\end{proposition}

\begin{beweis}[Beweisskizze]
Modulo~2: $a^2 \equiv a \pmod{2}$ für jedes ganze $a$. Also ist die Parität von
$a^2+b^2+c^2+d^2$ gleich der Parität der Summe $a+b+c+d$.
\begin{itemize}[noitemsep]
  \item Für $n \equiv 1 \pmod{4}$: Man rechnet modulo~4. Quadrate modulo~4 sind
        $0$ oder $1$. Eine Summe von vier solchen Werten ist $\equiv 1 \pmod{4}$
        genau dann, wenn genau eines oder drei der $a_i^2 \equiv 1 \pmod{4}$
        (d.\,h.\ genau 1 oder 3 ungerade Koeffizienten). Tatsächlich können auch
        alle vier ungerade sein (Summe $\equiv 0 \pmod{4}$, nein) --- eine sorgfältige
        Fallunterscheidung zeigt: es müssen 0, 2 oder 4 ungerade sein, wenn
        $n \equiv 1 \pmod{4}$, da sonst die Summe $\equiv 2$ oder $3 \pmod{4}$ wäre.
  \item Für $n \equiv 3 \pmod{4}$: Analog; hier müssen 1 oder 3 Koeffizienten ungerade sein.
\end{itemize}
\end{beweis}

\subsection{Tabelle: Konkrete Werte und Zerlegungen}

\begin{center}
\renewcommand{\arraystretch}{1.4}
\begin{tabular}{rccrrp{3.5cm}rr}
\toprule
$n$ & EABC & $\vii(3n{+}1)$ & $(a,b,c,d)$ & $a$ & Parität & $9n{+}6a{+}1$ & $\div 2^{\vii}$ \\
\midrule
 1 & E & 2 & $(1,0,0,0)$    &  1 & 1 ung., 3 ger. & 16  & 4 \\
 5 & E & 2 & $(2,1,0,0)$    &  2 & 2 ung., 2 ger. & 58  & 14 \\
   &   &   & $(1,2,0,0)$    &  1 & & 52  & 13 \\
 9 & C & 3 & $(3,0,0,0)$    &  3 & 1 ung., 3 ger. & 100 & 100/8=12,5? \\
   &   &   & \multicolumn{5}{l}{NB: $3\cdot 9+1=28$, $\nu_2(28)=2$} & 7 \\
 3 & B & 1 & $(1,1,1,0)$    &  1 & 3 ung., 1 ger. & 28  & 14 \\
 7 & A & 1 & $(2,1,1,1)$    &  2 & 3 ung., 1 ger. & 64  & 32 \\
   &   &   & $(1,2,1,1)$    &  1 & & 58  & 29 \\
11 & A & 1 & $(3,1,1,0)$    &  3 & 2 ung., 2 ger. & 118 & 59 \\
   &   &   & $(1,3,1,0)$    &  1 & & 106 & 53 \\
13 & C & 3 & $(3,2,0,0)$    &  3 & 1 ung., 3 ger. & 136 & 17 \\
   &   &   & $(2,3,0,0)$    &  2 & & 130 & --- \\
15 & B & 1 & $(3,2,1,1)$    &  3 & 3 ung., 1 ger. & 154 & 77 \\
21 & C & 3 & $(4,2,1,0)$    &  4 & 2 ung., 2 ger. & 214 & 214/8? \\
   &   &   & \multicolumn{5}{l}{NB: $3\cdot 21+1=64$, $\nu_2(64)=6$} & \\
\bottomrule
\end{tabular}
\end{center}

\begin{tcolorbox}[warnung]
\textbf{Beobachtung aus der Tabelle.}
Die Wahl von $a$ (Realteil) in der Vier-Quadrate-Zerlegung hängt von der
Zerlegungsstrategie ab. Für jede Zerlegung erhält man ein anderes $N(3q+1) = 9n+6a+1$.
Ob $9n+6a+1$ durch eine hohe Zweierpotenz teilbar ist (hohe $\vii$-Werte), hängt
von der Kongruenzklasse von $n$ ab --- nicht von der spezifischen Wahl von $a$.

Das bedeutet: Die EABC-Klasse ist eine Eigenschaft von $n$ modulo kleine Potenzen
von~2, unabhängig von der Quaternionendarstellung.
\end{tcolorbox}

\begin{proposition}[Unabhängigkeit von EABC-Klasse und Quaternionenwahl]
\label{prop:unabhaengigkeit}
Die 2-adische Valuation $\vii(3n+1)$ ist eine Eigenschaft von $n$ allein
(insbesondere von $n$ modulo~$4$ bzw.\ modulo~$8$ usw.) und hängt \emph{nicht}
von der Wahl der Vier-Quadrate-Zerlegung $n = a^2+b^2+c^2+d^2$ ab.
\end{proposition}

\begin{beweis}
$3n+1$ hängt nur von $n$ ab, nicht von der Darstellung von $n$ als Quaternionennorm.
$\vii(3n+1)$ ist eine Funktion von $3n+1$ allein.
\end{beweis}

\begin{bemerkung}
Proposition~\ref{prop:unabhaengigkeit} klärt die Grenzen des Quaternionen-Ansatzes:
Die EABC-Klasse ist keine Eigenschaft der Quaternionendarstellung, sondern der Zahl $n$
selbst. Die Quaternionendarstellung bereichert die Analyse durch die Formel
$N(3q+1) = 9n+6a+1$, aber die Frage, \emph{wie weit} $3n+1$ durch~2 teilbar ist,
bleibt eine modulo-arithmetische Frage.
\end{bemerkung}

% ====================================================
\section{Minimierungsproblem: Optimale Quaternionenwahl}
\label{sec:minimierung}
% ====================================================

\subsection{Problemformulierung}

\begin{definition}[Optimale Quaternionenwahl]
\label{def:optimal}
Für gegebenes $n \in \N$ definieren wir das \emph{Minimierungsproblem der Vier-Teilnormen}:
\[
  a_{\min}(n) \;:=\; \min\bigl\{ a \in \Z \;\big|\; \exists\, b,c,d \in \Z:\;
    a^2+b^2+c^2+d^2 = n \bigr\}.
\]
Das entsprechende Quaternion $q_{\min} = a_{\min}(n) + bi + cj + dk$ minimiert
$N(3q+1) = 9n + 6a + 1$ unter allen Zerlegungen von $n$.
\end{definition}

\begin{proposition}[Schranken für $a_{\min}$]
\label{prop:amin}
Für alle $n \geq 1$ gilt:
\[
  -\lfloor\sqrt{n}\rfloor \;\leq\; a_{\min}(n) \;\leq\; 0,
\]
mit Gleichheit $a_{\min}(n) = -\lfloor\sqrt{n}\rfloor$ genau dann, wenn
$n - \lfloor\sqrt{n}\rfloor^2$ als Summe von drei Quadraten darstellbar ist
(was nach dem Satz von Legendre für fast alle $n$ gilt).
\end{proposition}

\begin{beweis}
Da $a^2 \leq n$, gilt $|a| \leq \lfloor\sqrt{n}\rfloor$, also $a \geq -\lfloor\sqrt{n}\rfloor$.
Die Wahl $a = -\lfloor\sqrt{n}\rfloor$ ist genau dann realisierbar, wenn
$n - \lfloor\sqrt{n}\rfloor^2 = b^2+c^2+d^2$ für ganzzahlige $b,c,d$.
Nach dem Satz von Legendre (Dreieckszahlensatz) ist $m = b^2+c^2+d^2$ genau
dann möglich, wenn $m$ nicht von der Form $4^a(8b+7)$ ist.
\end{beweis}

\subsection{Minimales \texorpdfstring{$N(3q+1)$}{N(3q+1)}}

\begin{proposition}[Minimale Transformation]
\label{prop:minimal}
Das minimale Bild unter der Collatz-Quaternionen-Abbildung ist:
\[
  N_{\min}(n) \;:=\; \min_{q:\,N(q)=n} N(3q+1)
             \;=\; 9n + 6\,a_{\min}(n) + 1
             \;\leq\; 9n - 6\lfloor\sqrt{n}\rfloor + 1.
\]
\end{proposition}

\begin{korollar}[Grobe Abschätzung]
Für alle $n \geq 1$ gilt $N_{\min}(n) \leq 9n - 6\lfloor\sqrt{n}\rfloor + 1$.
Der Term $-6\lfloor\sqrt{n}\rfloor$ liefert eine Kontraktion von $O(\sqrt{n})$
gegenüber dem isotropen Bild $9n+1$.
\end{korollar}

\begin{bemerkung}
Für große $n$ gilt $N_{\min}(n) \approx 9n - 6\sqrt{n}$, was immer noch
super-linear in $n$ wächst. Die Kontraktion durch negative Wahl von $a$
reicht nicht aus, um die Collatz-Dynamik allein durch Quaternionen-Optimierung
zu kontrollieren --- der entscheidende Faktor bleibt die Division durch
$2^{\vii(3n+1)}$ nach dem ungeraden Schritt.
\end{bemerkung}

\subsection{Spezialfall: Quadratzahlen}

\begin{theorem}[Quadratzahlen und Quaternionentransformation]
\label{thm:quadratzahlen}
Sei $n = m^2$ eine Quadratzahl mit $m \geq 1$. Mit der Wahl $q = -m$ (rein reell)
gilt $a = -m$, $b = c = d = 0$, $N(q) = m^2 = n$, und:
\[
  N(3q + 1) \;=\; (3(-m) + 1)^2 \;=\; (1 - 3m)^2 \;=\; (3m-1)^2.
\]
Das Bild $N(3q+1) = (3m-1)^2$ ist ebenfalls eine Quadratzahl!
\end{theorem}

\begin{beweis}
Direkte Rechnung: $3q+1 = 3(-m)+1 = 1-3m$, also
$N(3q+1) = (1-3m)^2 = (3m-1)^2$. \qedhere
\end{beweis}

\begin{tcolorbox}[ergebnis, title={Schöner Spezialfall: Quadratzahlen}]
Für $n = m^2$ gilt mit der optimalen Quaternionenwahl $q = -m$:
\[
  N(3q+1) \;=\; (3m-1)^2.
\]
Die Tabelle der ersten Werte:
\begin{center}
\renewcommand{\arraystretch}{1.3}
\begin{tabular}{rrrrr}
\toprule
$m$ & $n = m^2$ & $q = -m$ & $N(3q+1) = (3m-1)^2$ & $\sqrt{N(3q+1)}$ \\
\midrule
1 &  1 & $-1$ &   4 &  2 \\
2 &  4 & $-2$ &  25 &  5 \\
3 &  9 & $-3$ &  64 &  8 \\
4 & 16 & $-4$ & 121 & 11 \\
5 & 25 & $-5$ & 196 & 14 \\
6 & 36 & $-6$ & 289 & 17 \\
\bottomrule
\end{tabular}
\end{center}
Die Folge $1, 5, 8, 11, 14, 17, \ldots$ hat die Differenz~3: $3m-1 = 2, 5, 8, 11, \ldots$
für $m = 1, 2, 3, \ldots$ --- eine arithmetische Progression mit Differenz~3.
\end{tcolorbox}

\begin{bemerkung}[Verbindung zum Collatz-Schritt für Quadratzahlen]
Für $n = m^2$ ist der Collatz-Schritt $3n+1 = 3m^2+1$. Die Zerlegungsformel
$N(3q+1) = (3m-1)^2$ beschreibt jedoch nicht den Collatz-Schritt $3n+1$,
sondern die Norm des Bildes des Quaternions unter der Abbildung $q \mapsto 3q+1$
bei optimaler Wahl $q = -m$.

Die entscheidende Frage: Wird $(3m-1)^2$ als Norm des Collatz-Bildes realisiert?
Das Collatz-Bild von $n = m^2$ ist $3m^2+1$ (für ungerades $m$), und
$(3m-1)^2 = 9m^2 - 6m + 1 \neq 3m^2 + 1$ im Allgemeinen.

Der Quaternionen-Ansatz kodiert also eine andere (optimierte) Trajektorie als
die Standard-Collatz-Abbildung.
\end{bemerkung}

% ====================================================
\section{Verbindung zur D\texorpdfstring{${}_4$}{4}-Theta-Funktion und \texorpdfstring{$\mathcal{I}(n)$}{I(n)}}
\label{sec:theta}
% ====================================================

\subsection{Die D${}_4$-Theta-Funktion}

Das D${}_4$-Gitter (Hurwitz-Quaternionen-Gitter) hat die Theta-Funktion:
\[
  \Theta_{D_4}(\tau) \;=\; \sum_{q \in \mathcal{H}} e^{i\pi\tau N(q)}
  \;=\; \sum_{n=0}^{\infty} r_4^{\mathcal{H}}(n)\, q^n,
  \quad q = e^{2\pi i \tau},
\]
wobei $r_4^{\mathcal{H}}(n)$ die Anzahl der Hurwitz-Quaternionen mit Norm $n$ ist.

Die Theta-Funktion des D${}_4$-Gitters ist bekannt als:
\[
  \Theta_{D_4}(\tau) \;=\; \frac{1}{2}\left(\theta_3(\tau)^4 + \theta_2(\tau)^4 + \theta_4(\tau)^4 - 2\right),
\]
wobei $\theta_2, \theta_3, \theta_4$ die Jacobi-Theta-Funktionen sind.

Alternativ gilt für das einfachere Lipschitz-Gitter:
\[
  \Theta_{\mathcal{L}}(\tau) \;=\; \theta_3(\tau)^4
  \;=\; \left(\sum_{n \in \Z} q^{n^2}\right)^4
  \;=\; \sum_{n=0}^{\infty} r_4(n)\, q^n,
\]
was den Jacobi-Vierquadratensatz reproduziert.

\subsection{Verbindung zu Eisenstein-Reihen und Bernoulli-Zahlen}

Die Eisenstein-Reihen $G_{2k}(\tau) = 2\zeta(2k) + \frac{2(2\pi i)^{2k}}{(2k-1)!}
\sum_{n=1}^\infty \sigma_{2k-1}(n) q^n$ sind Modulformen, die durch Gitterpunktsummen
definiert sind. Die Verbindung zur D${}_4$-Theta-Funktion:

\begin{proposition}[Eisenstein-Reihen und Vier-Quadrate]
\label{prop:eisenstein}
Die Theta-Funktion $\Theta_{\mathcal{L}}(\tau) = \theta_3(\tau)^4$ ist eine
Modulform vom Gewicht~2 für eine Kongruenzuntergruppe von $\mathrm{SL}_2(\Z)$.
Für die ersten Koeffizienten gilt:
\[
  r_4(n) \;=\; 8 \sum_{\substack{d \mid n \\ 4 \nmid d}} d
          \;=\; 8\,\sigma_1(n) - 32\,\sigma_1(n/4),
\]
wobei $\sigma_1(m) = \sum_{d|m} d$ die Summe der Teiler ist und $\sigma_1(m) = 0$
für nicht-ganzzahliges $m$.
\end{proposition}

\subsection{Hypothetische Verbindung zu \texorpdfstring{$\In(n)$}{I(n)}}

\begin{tcolorbox}[lücke, title={Offene Frage: D${}_4$-Theta und Bernoulli-Normschalen}]
\textbf{Hypothese (unbewiesen):} Gibt es eine direkte Verbindung zwischen
$r_4(n)$ (Anzahl der Vier-Quadrate-Zerlegungen) und $\In(n)$ (Bernoulli-Normschale)?

Bekannte Verbindungen:
\begin{enumerate}[leftmargin=*, noitemsep, label=(\alph*)]
  \item Sowohl $r_4(n)$ als auch $\In(n)$ sind mit Eisenstein-Reihen verknüpft.
  \item Die Formel $\In(n) = \frac{2(2n)!\,\zeta(2n)}{(2\pi)^{2n}}\cdot D(n)$
        verwendet $\zeta(2n)$, das seinerseits via Modulformen mit Theta-Funktionen
        zusammenhängt.
  \item Beide Größen messen auf unterschiedliche Weise die ``Dichte'' von
        Gitterpunkten --- $r_4(n)$ auf dem Lipschitz-Gitter, $D(n)$ auf dem
        Primzahlgitter (von Staudt--Clausen).
\end{enumerate}

Eine explizite Formel, die $r_4(n)$ und $\In(n)$ verbindet, ist offen.
\end{tcolorbox}

\begin{bemerkung}[Modulare Struktur]
Die Theta-Funktion $\theta_3(\tau)^4$ und die Eisenstein-Reihe $E_2(\tau)$ sind
beide Modulformen (bzw.\ quasi-Modulformen) für $\Gamma_0(4)$. Die Bernoulli-Zahlen
erscheinen als Konstanten in den $q$-Entwicklungen der Eisenstein-Reihen.
Die Verbindung zwischen $r_4(n)$ und $B_{2n}$ verläuft daher durch den
Raum der Modulformen --- ein tiefer, aber noch nicht vollständig verstandener Link.
\end{bemerkung}

% ====================================================
\section{Fazit}
\label{sec:fazit}
% ====================================================

\subsection{Zusammenfassung der Ergebnisse}

\begin{center}
\renewcommand{\arraystretch}{1.5}
\begin{tabular}{p{7cm} p{2.8cm} p{3.5cm}}
\toprule
\textbf{Aussage} & \textbf{Status} & \textbf{Anmerkung} \\
\midrule
$N(3q+1) = 9N(q) + 6\,\Re(q) + 1$
  & \textcolor{bewiesengrün}{\bfseries Beweis}
  & Direktrechnung \\
Transformation der Teilnormen (Gleichungen \ref{eq:N1}--\ref{eq:N4})
  & \textcolor{bewiesengrün}{\bfseries Beweis}
  & Direktrechnung \\
Für $n = m^2$: $N_{\min}(3q+1) = (3m-1)^2$
  & \textcolor{bewiesengrün}{\bfseries Beweis}
  & Theorem~\ref{thm:quadratzahlen} \\
EABC-Klasse unabhängig von Quaternionenwahl
  & \textcolor{bewiesengrün}{\bfseries Beweis}
  & Proposition~\ref{prop:unabhaengigkeit} \\
Paritätsmuster nach Residuenklasse mod~4
  & \textcolor{bewiesengrün}{\bfseries Beweis}
  & Proposition~\ref{prop:paritaet} \\
\midrule
Verbindung D${}_4$-Theta zu $\In(n)$
  & \textcolor{warnorange}{\bfseries Offen}
  & Strukturell motiviert \\
Optimale Quaternionenwahl steuert Konvergenz
  & \textcolor{warnorange}{\bfseries Heuristik}
  & Kontraktion zu gering \\
\midrule
Quaternionenwahl liefert direkten Collatz-Beweis
  & \textcolor{lückenrot}{\bfseries Falsch}
  & EABC-Klasse unabhängig \\
\bottomrule
\end{tabular}
\end{center}

\subsection{Strukturelle Erkenntnisse}

\begin{tcolorbox}[hauptsatz]
\textbf{Hauptergebnisse:}

\medskip
\textbf{(1) Asymmetrische Transformation.}
Der Collatz-Schritt $q \mapsto 3q+1$ auf $\mathcal{L}$ transformiert die
vier Teilnormen asymmetrisch:
\[
  (N_1, N_2, N_3, N_4) \;\mapsto\;
  \bigl(9N_1 + 6\sqrt{N_1}+1,\; 9N_2,\; 9N_3,\; 9N_4\bigr).
\]
Die Realteil-Richtung ist die \textbf{ausgezeichnete Richtung}, die den
gesamten Collatz-Korrekturfaktor $6a+1$ trägt.

\medskip
\textbf{(2) Quadratzahlformel.}
Für $n = m^2$ mit optimaler Wahl $q = -m$ gilt $N(3q+1) = (3m-1)^2$ ---
eine Quadratzahl. Die Folge $3m-1$ ist eine arithmetische Progression.

\medskip
\textbf{(3) Grenzen des Ansatzes.}
Die EABC-Klasse (2-adische Valuation von $3n+1$) ist eine Eigenschaft von $n$
allein und hängt nicht von der Quaternionendarstellung ab.
Die Quaternionen-Zerlegung bereichert die Struktur, kontrolliert aber nicht
die 2-adische Tiefe des Collatz-Schritts.

\medskip
\textbf{(4) Ausblick: D${}_4$-Theta.}
Die Vier-Quadrate-Zerlegung ist über die D${}_4$-Theta-Funktion mit den
Eisenstein-Reihen und (über Euler) mit den Bernoulli-Zahlen verknüpft.
Eine explizite Verbindung zu den Bernoulli-Normschalen $\In(n)$ ist offen
und strukturell motiviert.
\end{tcolorbox}

\subsection{Offene Fragen}

\begin{tcolorbox}[lücke]
\begin{enumerate}[leftmargin=*, label=(\arabic*)]
  \item \textbf{Modulare Verbindung $r_4(n)$ und $\In(n)$:}
        Gibt es eine Formel, die die Anzahl der Vier-Quadrate-Zerlegungen
        mit dem Bernoulli-Zähler $\In(n)$ verbindet?

  \item \textbf{Hurwitz-Quaternionen und EABC:}
        Haben Primzahlen $p \equiv 1 \pmod 4$ (E-Klasse) strukturell andere
        Hurwitz-Quaternionen-Darstellungen als $p \equiv 3 \pmod 4$ (A/B-Klasse)?

  \item \textbf{Optimierungsproblem in höheren Dimensionen:}
        Gibt es eine Verallgemeinerung auf Oktonionen (8 Teilnormen), die eine
        feinere Struktur der Collatz-Abbildung aufdeckt?

  \item \textbf{Spektrale Interpretation:}
        Das D${}_4$-Gitter hat eine natürliche spektrale Interpretation via
        Theta-Reihen. Kann die Collatz-Dynamik als Fluss auf dem D${}_4$-Gitter
        verstanden werden?
\end{enumerate}
\end{tcolorbox}

% ====================================================
% Literatur
% ====================================================
\begin{thebibliography}{99}

\bibitem{hoffbauer-bernoulli}
T.\ Hoffbauer,
\textit{Bernoulli-Normschalen und Collatz-Dynamik: Der Satz von von Staudt--Clausen
als Lyapunov-Kandidat},
Mathematische Texte, \texttt{collatz\_bernoulli\_schalen.tex} (2026).

\bibitem{hoffbauer-lyapunov}
T.\ Hoffbauer,
\textit{Der Lyapunov-Exponent der Collatz-Dynamik},
Mathematische Texte, \texttt{collatz\_lyapunov.tex} (2026).

\bibitem{hoffbauer-oktonionen}
T.\ Hoffbauer,
\textit{Collatz-Beweis via Oktonionen-Norm},
Mathematische Texte, \texttt{collatz\_oktonionen\_beweis.tex} (2026).

\bibitem{lagrange}
J.-L.\ Lagrange,
\textit{Démonstration d'un théorème d'arithmétique},
Nouv.\ Mém.\ Acad.\ Berlin (1770), 123--133.

\bibitem{jacobi}
C.\,G.\,J.\ Jacobi,
\textit{Fundamenta Nova Theoriae Functionum Ellipticarum},
Königsberg (1829).

\bibitem{hurwitz}
A.\ Hurwitz,
\textit{Vorlesungen über die Zahlentheorie der Quaternionen},
Springer, Berlin (1919).

\bibitem{conway-sloane}
J.\,H.\ Conway, N.\,J.\,A.\ Sloane,
\textit{Sphere Packings, Lattices and Groups},
3.\ Aufl., Springer, New York (1999).

\bibitem{serre-cours}
J.-P.\ Serre,
\textit{Cours d'arithmétique},
Presses Universitaires de France, Paris (1970).

\bibitem{ireland-rosen}
K.\ Ireland, M.\ Rosen,
\textit{A Classical Introduction to Modern Number Theory},
Springer, New York, 2.\ Aufl.\ (1990).

\bibitem{tao-collatz}
T.\ Tao,
\textit{Almost all orbits of the Collatz map attain almost bounded values},
Forum of Mathematics Pi \textbf{10} (2022), e12.
\texttt{arXiv:1909.03562}

\end{thebibliography}

\end{document}
